\section{Exact reuse of Bellman certificate geometry}
\label{supp:cachedgeometry126c}

The following results concern the separately frozen within-service cache.
They apply to a finite nonempty collection of rational observations with
$A>0$ and $0<\mu\le L$.  Every observation anchor, requested action and integration endpoint
is feasible.  A case and its payload entries are immutable throughout its
lifetime inside one service.  The normalized cases have a fixed number of scalar fields outside their
finite payload list.  The implementation retains a strong reference
to every case used for an object-identity binding; thus a discarded temporary
case cannot have its identifier reassigned to a different case while the
binding exists.  Mutating a bound case is outside these assumptions.

Let $s$ be the number of retained derivative observations, including repeated
anchors.  The two transported functions of observation $j$ are
\begin{align*}
 \lambda_j(t)&=g_j^-+
  \begin{cases}L(t-a_j),&t<a_j,\\\mu(t-a_j),&t\ge a_j,\end{cases}\\
 \upsilon_j(t)&=g_j^++
  \begin{cases}\mu(t-a_j),&t<a_j,\\L(t-a_j),&t\ge a_j.\end{cases}
\end{align*}
Write $e_-(t)=\max_j\lambda_j(t)$ and
$e_+(t)=\min_j\upsilon_j(t)$.  These are precisely the inherited
transported envelopes, without adding an unrecorded information constraint.
Their geometry key is
\[
 \kappa_D=(A,\mu,L,((a_j,g_j^-,g_j^+))_{j=1}^s).
\]
The observation order is retained.  A subset request uses the corresponding
ordered subset in this key.  The inherited algorithm first partitions the
domain at all observation anchors.  On each resulting open interval it
retains every strict interior intersection of nonparallel transported lines,
even when those lines do not form the active envelope.  On each subinterval
it chooses the extremizing line at the midpoint, with the original fixed
lexicographic rule for ties.  No segments are coalesced by the cache.

For either envelope, write its full ordered partition as
\[
 (x_i,x_{i+1},m_i,c_i)_{i=0}^{P-1},\qquad
 0=x_0<\cdots<x_P=A,
 \qquad e(t)=m_it+c_i\quad(x_i<t<x_{i+1}).
\]
At endpoints use the continuous extension.  Define exact prefixes
\[
 I_0=0,\qquad
 I_{i+1}=I_i+\frac{m_i}{2}(x_{i+1}^2-x_i^2)
                              +c_i(x_{i+1}-x_i).
\]
For $t\in[x_i,x_{i+1}]$ the stored primitive evaluates
\begin{equation}\label{eq:cacheprimitive126c}
 F(t)=I_i+\frac{m_i}{2}(t^2-x_i^2)+c_i(t-x_i).
\end{equation}
The values at $0$ and $A$ are explicitly initialized.  Binary search on the
ordered starts chooses the interval for every other new endpoint; a boundary
can be assigned to either adjacent interval because the two formulas agree.

\begin{proposition}[Exact geometry and service identity]
\label{prop:cacheidentity126c}
Under the stated assumptions, the cache returns exactly the same rational
result as the inherited implementation for each of the following operations:
full or partial envelope partition; both directed contrast endpoints, including
the best-single-record mode; derivative-envelope value; constrained envelope
zero; and raw-support minimum with its selected minimizing action.

The raw-support minimum is keyed separately by
\[
 \kappa_S=(A,\mu,((a_j,W_j,p_j^-,p_j^+))_{j=1}^s).
\]
The original incumbent, future allowance, original-optimum floor and upper
value levels are not arguments of that raw minimum.  They remain arguments
of the separately recalculated policy certificate.  In particular, a cached
raw minimum is not a cached value of the final original-optimum lower bound.

Suppose the original oracle and actor return the same outputs on the same
inputs, and the frozen service's control decisions depend only on scientific
values and declared discrete budgets, not on elapsed time.  Replacing the
listed geometry operations by their cached versions then leaves the complete
sequence of economic calls, candidate decisions and scientific certificates
unchanged whenever both executions complete their required arithmetic.  The
new cache counters and clocks may differ.  Resetting the cache at each paid
decision preserves this conclusion without transferring acquired information
between services.
\end{proposition}

\begin{proof}
First consider a partial interval $[u,v]\subseteq[0,A]$ with $u<v$.
Its old partition has the endpoints $u,v$, every observation anchor in its
interior, and every pairwise line crossing in the interior of each resulting
anchor interval.  These are exactly the full-domain breakpoints strictly
between $u$ and $v$, together with $u,v$.  To see this, restrict any full
anchor interval to $[u,v]$.  Each observation still has the same transported
line on that restriction.  Its intersection with any other nonparallel line
is the same rational number; it is retained precisely when it lies inside
the restriction.  A crossing at an anchor or at $u,v$ is already an endpoint.
Parallel equal lines create no extra crossing in either calculation.

The relative order of distinct affine lines is constant on each open
subinterval between all crossings.  Thus its active envelope line cannot
change when a subinterval is truncated at $u$ or $v$.  If several lines
agree there, their slope and intercept agree throughout that open interval,
so the fixed tie rule also gives the same coefficients.  Clipping the full
list therefore gives exactly the old partial list, including the retained
inactive-line crossings and coefficient pairs.  If $u=v$, both routines
return the empty list.  The same argument applies to a declared subset.

Induction on $i$ shows $I_i=\int_0^{x_i}e(t)\,dt$.  Integrating the
selected affine line on the final partial segment proves
\eqref{eq:cacheprimitive126c}.  Consequently its differences are the exact
integrals returned by summing the old partial partition.  Rational addition
introduces no order-dependent rounding.  For $a<b$, the lower contrast uses
$F_-(b)-F_-(a)$ and the upper one uses $F_+(b)-F_+(a)$.  For $a>b$,
reversing the integration direction reverses the two bounds.  For $a=b$,
both routines return zero.  The best-single-record mode applies this identity
to each singleton record and then takes the same maximum of lower endpoints
and minimum of upper endpoints.

Each envelope is continuous, because it is a finite maximum or minimum of
continuous transported functions.  Evaluating the selected affine piece at
an interior point, a crossing, an anchor or either domain endpoint therefore
gives the same value as direct evaluation of every observation.  All slopes
are at least $\mu>0$, so the envelope is strictly increasing.  The two
endpoint sign tests and, when needed, the first feasible zero $-c_i/m_i$
therefore return its identical constrained zero.  Saving an already computed
value or zero cannot change it.

The raw-support routine minimizes
\[
 S(t)=\max_j\left\{W_j+
        \min[p_j^-(t-a_j),p_j^+(t-a_j)]
                              +\frac\mu2(t-a_j)^2\right\}
 \quad(0\le t\le A).
\]
Every argument of this function is contained in $\kappa_S$; the omitted
fields do not enter it.  The cache executes the unchanged exact routine on
the first occurrence of this complete key, retains its value and original
tie-selected action, and returns that same pair thereafter.  Equality of
$\kappa_D$ likewise implies equality of every argument of the transported
geometry.  The object-identity binding is sound by immutability and its
strong references.  This proves all individual identities.

For the service conclusion, induct on the ordered external calls and
scientific control decisions.  At the first call the frozen inputs agree.
An original economic or actor call with the same inputs has the same output
by the stated premise.  A substituted geometry call has the same output by
the identities just proved.  A scientific branch or discrete-budget update
therefore takes the same branch in both executions.  The induction preserves
all subsequent oracle arguments, certificates, candidate refusals, accepted
changes and stopping decisions.  Cache-table accesses are certificate work;
they neither call an economic primitive nor replace a paid economic record.
Each complete service starts with an empty cache and the context restores the
original functions and clears its bindings on exit, including exceptional
exit.  A failure caused by unavailable machine resources is not ruled out by
an algebraic identity, and a failed execution is not asserted to complete.
\end{proof}

\begin{lemma}[Size and construction cost of the inherited partitions]
\label{lem:cachepieces126c}
For one sign of a derivative ledger with $s\ge1$ observations, the number
of retained full-domain pieces obeys
\begin{equation}\label{eq:cachepieces126c}
 P\le 1+s+\binom{s}{2}.
\end{equation}
The inherited all-crossing construction, followed by its prefix calculation,
uses $O(s^3)$ exact scalar arithmetic, ordering and item-inspection operations
and stores $O(s^2)$ piece data.  These are operation counts, not bounds on
the bit cost of rational numerators and denominators.

A first exact raw-support minimum for a key with $s$ supports also has an
$O(s^3)$ bound in this operation model.  It need not be a convex function,
and its minimum is not obtained by assuming convexity of its slope-interval
envelope.
\end{lemma}

\begin{proof}
There are at most $s$ distinct internal observation anchors.  Consider two
anchors $a_j\le a_k$ and the difference of their lower transported functions.
It is constant on each exterior region $t<a_j$ and $t>a_k$, because both
functions have the same slope there.  On $(a_j,a_k)$ its slope is
$\mu-L\le0$.  For the upper transported functions the middle slope is
$L-\mu\ge0$.  Thus a pair of either kind has at most one isolated interior
crossing not already an anchor.  If $L=\mu$ it has none.  Coincident lines
on an interval contribute no isolated crossing.  Counting all pairs proves
\eqref{eq:cachepieces126c}, even though the algorithm retains crossings of
inactive lines and repeated anchors.

There are at most $s+1$ anchor intervals.  Forming their $s$ lines and
checking every pair costs $O(s^3)$ operations.  Let $P$ be the actual
number of pieces.  Sorting the retained crossings costs at most
$O((P+s)\log(s+2))$, and selecting the extremizing line on each piece
costs $O(sP)$.  With \eqref{eq:cachepieces126c}, both are $O(s^3)$.
The prefixes add $O(P)$ operations and storage.  These statements describe
the inherited enumeration; no optimal lower-envelope construction algorithm
is being substituted in the complexity calculation.

For the support minimum subtract the common quadratic $\mu t^2/2$.
Each remaining support is a continuous piecewise-affine function with one
break at its own anchor.  A pair has at most three affine regions, hence
at most three isolated crossings.  There are consequently at most
$1+s+3\binom{s}{2}$ resulting pieces before selecting their maximum.
The inherited routine enumerates all pairwise crossings on every anchor
interval and selects the maximal line on each piece.  It retains all
piece endpoints and the feasible stationary point of the common quadratic
plus its selected affine line.  These candidates include a minimizer on
every closed piece.  Evaluating the original support maximum at the
candidates and choosing its minimum requires $O(s)$ work per candidate.
There are $O(s^2)$ candidates and $O(s^3)$ pair-enumeration work, yielding
the asserted bound.  Downward slope-interval kinks cause no gap in this
piecewise argument.
\end{proof}

\begin{proposition}[Finite operation and storage account for geometry reuse]
\label{prop:cachework126c}
Within one paid service let $\mathcal D$ be the distinct derivative keys and
$\mathcal S$ the distinct raw-support keys.  Write $s_\kappa$ for the
support count at a key and $P_\kappa$ for the total number of its derivative
pieces actually built across both signs.  For a derivative key let
$e_\kappa$ count previously unseen primitive arguments,
$v_\kappa$ previously unseen gradient arguments, and $l_\kappa$
nondegenerate requests for a clipped partition.  Let $Q$ count all requests
to the cache routines, including calls made by another cache routine, and
let $C$ count the total clipped pieces returned.  Let $B$ count formations
of a derivative key, and $R$ count requests for a raw-support minimum;
write $s_b,s_r$ for their respective record lengths.

There are at most $2|\mathcal D|$ full derivative-partition builds,
$2|\mathcal D|$ initial constrained-root calculations, and
$|\mathcal S|$ executions of the original raw-support-minimum routine.
The explicit geometry arithmetic, comparisons and key-entry inspections
are bounded by
\begin{align}
 O\Bigg(&\sum_{\kappa\in\mathcal D}
   \left[s_\kappa^3+P_\kappa+
    (e_\kappa+v_\kappa+l_\kappa)\log(2+P_\kappa)\right]
   +\sum_{\sigma\in\mathcal S}s_\sigma^3 \nonumber\\
 &\hspace{12mm}+Q+C+\sum_{b=1}^{B}s_b+\sum_{r=1}^{R}s_r\Bigg).
 \label{eq:cachework126c}
\end{align}
In particular, if every record has at most $S$ entries and each built
geometry at most $P$ pieces, this is
\[
 O\big((|\mathcal D|+|\mathcal S|)S^3
                   +Q[S+\log(2+P)]+C\big).
\]
The routines also make $O(Q)$ cache-table operations.  This counts the
operations requested of the tables and does not assume a constant bound
on their internal hashing, collision handling or allocation time.  Rational
bit arithmetic and those table costs are additional machine work, included
in the measured complete service.

Cached geometry, memoized arguments, exact keys and retained immutable
case references require at most
$O(\sum_{\kappa\in\mathcal D}P_\kappa+QS)$ stored scalar or reference
entries under the same bound $S$.  This memory is confined to one service.
The result establishes neither a lower bound for the uncached computation
nor a uniform complete-time advantage for the cached one.
\end{proposition}

\begin{proof}
Each derivative key has two possible sign entries and each entry is built
only on its first request.  Lemma~\ref{lem:cachepieces126c} pays for those
builds and the prefix sums.  The first gradient or primitive evaluation at
an argument uses binary search on ordered rational starts, requiring
$O(\log(2+P_\kappa))$ comparisons, followed by a constant number of
exact scalar operations.  Repeated arguments retrieve the exact saved
value.  A clipped partition pays for a binary search and one iteration per
returned piece.  Each initial root scans at most its stored partition;
there are only two roots per key, so these scans are covered by the
$P_\kappa$ term.  The four primitive evaluations and endpoint reversal
in a fused contrast are included among the requests.  Recursive
best-single contrasts also enter the request count rather than being
charged as one constant-cost operation.

The raw-support routine runs once per distinct complete support key.
Importantly, its current implementation forms that key on every request,
including a cache hit.  Its $s_r$ entry inspections therefore remain in
\eqref{eq:cachework126c}; they are not replaced by an unproved constant
cost.  A previously bound immutable case avoids rebuilding its derivative
key, whereas each fresh binding or explicit subset forms all $s_b$ entries.
Every request performs only a fixed number of additional table operations,
apart from recursive requests already counted in $Q$.  Summing these costs
proves the displayed bound.  Since key formations and minimum requests are
included among the calls, replacing their lengths by $S$ gives the simpler
bound.

Each built piece stores a fixed number of rational entries and a prefix.
The number of memoized arguments, roots, bindings and support keys is at
most a constant multiple of the request count.  A mathematical key or a
retained case contains at most $O(S)$ scalar or reference entries.  These
facts prove the storage bound.  Clearing the caches on context exit prevents
these retained cases or evaluations from becoming an uncharged initial state
for the next service.

The upper bound counts the stated calculations.  An uncached partial
interval can contain fewer pieces than the full domain, or a service may
never request the same geometry twice.  Its work therefore has no matching
lower bound merely because the cached algorithm has the displayed upper
bound.  Growing rational bit lengths, table operations, extra retained
objects and cache setup or cleanup can also affect actual cost.  The paired
experiment measures all those costs while requiring identical original
economic work and policies.  It supplies any empirical speed comparison;
the algebra alone supplies none.
\end{proof}
